\chapter{The competition axis: relational collateral form}
\label{ch:competition}

\section{Why ability is relational}

The hardest and most under-exploited dimension of the problem is \emph{competition}. The handicap
system actively compresses the field --- the handicapper assigns weight to engineer a blanket
finish --- which is precisely why win probabilities are close and why beating the market is hard.
The consequence for modelling is that \emph{raw} ability is the wrong object: beating a strong
field should count for more than the same performance against a weak one. The win model, at
equilibrium, \emph{is itself a rating system}; the right ability measure is the consistent
solution of an opponent-adjusted system.

The first attempt at this axis --- a scalar Elo / relative-form rating --- failed, and the failure
is instructive. A scalar rating is collinear ($\approx 0.70$) with the official rating the market
already prices, and it discards the only thing competition actually does, which is to make outcomes
\emph{relational}: a horse's effective ability is a \emph{function of the opponent set}, not a
context-free number. A front-runner among closers steals an uncontested lead and arrives with
reserve unspent; the same front-runner among other front-runners is forced into a contested duel
that burns $\Wp$ against the air with no draft, and fades. Identical horse, different field,
different result. A scalar cannot represent a quantity that depends on who else is in the race. The
correct object is \emph{relational}, and the relational treatment converts strongly (HV $+2.3$pp,
with ECE improving) where the scalar one vanished --- the canonical illustration that ``a dead
first attempt is not a dead axis.''

\section{The Massey collateral-form rating}

\subsection{Construction}

The collateral-form rating is an opponent-adjusted \emph{velocity} rating in m/s, built as an
online Massey least-squares system so that it is strictly leak-free. The atom is again the
condition-normalised speed figure of a run, $s_i = L/\Tfin_i - \mathrm{par}(\text{going},L)$ (m/s,
the same kinematic figure as \eqref{eq:speedfig}). Races are processed in date order. The
\emph{field strength} of a race $R$ is the mean \emph{prior} collateral rating $\bar R$ of its
runners (unrated newcomers count as neutral $0$). A runner's observation from race $R$ combines its
\emph{margin against the field it actually ran} with the \emph{strength of that field}:
\begin{equation}
  \mathrm{obs}_i \;=\; \underbrace{\bigl(s_i - \overline{s}_R\bigr)}_{\text{margin vs this field}}
   \;+\; \kappa\,\underbrace{\bar R}_{\text{strength of this field}},
  \label{eq:collat-obs}
\end{equation}
and the horse's rating is the expanding mean of its $\mathrm{obs}$ over its strictly-prior runs.
The $\kappa\bar R$ term is the \emph{collateral propagation}: a margin posted against a strong
field inherits that field's strength, which is one Gauss--Seidel sweep of the Massey
least-squares system in date order. The least-squares value is $\kappa=1$; $\kappa<1$ damps the
propagation. Leak-freeness is preserved by emitting the rating a horse carries \emph{into} a race
(the running per-horse mean over strictly-prior runs) and only \emph{then} updating that running
mean with the current race's observation --- the same strictly-prior guarantee as the form table,
and the par is a condition constant rather than the horse's own outcome, so the full panel is safe.

\subsection{Coupling}

The emitted $\mathrm{collat\_form}$ is within-race centred (a field-mean shift cancels in the win
probability) and coupled to the drift with a per-venue gain $\texttt{COLLAT\_GAIN\_BY\_VENUE}$ (ST
$1.0$, HV $2.5$). The much larger HV gain is consistent with HV's tighter, more compressed fields
making opponent adjustment more informative there; the collateral channel closes roughly a third
of the ST--HV skill gap.

\section{Deepening the rating (v2)}

The production ST collateral rating uses a refined construction (\code{collateral\_form\_table\_v2})
that pushes the Massey solution closer to its true fixed point and conditions it more carefully.
All refinements preserve the strictly-prior / online leak-free discipline. The baked ST
construction is $\{\,n_{\text{sweeps}}=3,\ \text{half\_life}=4.0,\ \text{band\_par}=\text{True}\,\}$;
HV keeps the single-sweep table (band conditioning fragments its tight trip range). The four
refinement axes:

\begin{enumerate}[label=\textbf{R\arabic*.},leftmargin=*]
  \item \textbf{Deeper propagation ($n_{\text{sweeps}}>1$).} A single online pass propagates
        opponent strength only one hop along the collateral chain. Additional \emph{Jacobi}
        fixed-point sweeps re-solve the strictly-prior race graph before each new race-date is
        emitted, each sweep recomputing every prior race's observation using the previous sweep's
        \emph{frozen} ratings as the field-strength source. Strength then propagates one more hop
        per sweep, toward the true Massey fixed point. Crucially the extra sweeps touch only
        strictly-prior races, so the rating emitted for a race still contains no information from
        it or any later race. Baked at three sweeps on ST.
  \item \textbf{Recency weighting (\emph{half\_life}).} The per-horse rating can be an
        exponentially-weighted mean of its observations rather than a flat expanding mean, so a
        stale opponent rating decays --- the recency-form lever (Section \ref{ch:features}) applied
        to the competition channel. With decay $\delta = 2^{-1/\text{half\_life}}$ the running
        numerator and denominator are updated as $\delta\cdot(\text{old}) + (\text{new})$. Baked at
        a four-race half-life on ST.
  \item \textbf{Nonlinear margin (\emph{margin\_pow}) and strength-of-schedule asymmetry
        (\emph{asym}).} The within-race margin can be raised to a sign-preserving power
        ($\mathrm{sign}(m)\,|m|^{p}$, $p>1$ rewards big margins, $p<1$ compresses), and the field
        strength inherited on a \emph{loss} can be weighted differently from a \emph{win} (kept
        term $\kappa\bar R$ for a beat, $\kappa\cdot\mathrm{asym}\cdot\bar R$ for a loss), encoding
        that losing to a strong field tells you less than beating one. Both default to the linear,
        symmetric baked values.
  \item \textbf{Trip-band conditioning (\emph{band\_par}).} The par velocity and the margin are
        computed within (going, distance-\emph{band}), and the field-strength inheritance is
        conditioned on the trip band, so a 1200\,m margin feeds a horse's rating only for nearby
        trips. Ability is partly trip-specific; band conditioning enforces like-for-like
        comparison. A horse carries a per-band rating with a fallback to its global mean. Baked on
        ST.
\end{enumerate}

The deepened ST construction added $+1.3$pp \emph{on top of} the trial fusion (super-additive),
and was what broke ST through $65\%$; its ECE also improves. This is the competition axis,
operationalised --- the relational rating the ``win model is a rating system at equilibrium''
framing predicts, built jointly with the leak-free discipline rather than bolted on.

\begin{remark}[The free theory and the interacting theory]
The production Wald/PL surrogate is the \emph{free}, non-interacting theory: context-free drifts,
no traffic, no draft. The collateral rating injects the \emph{first relational moment} of
competition --- opponent strength --- into the per-horse drift, which is why it survives the
collinearity that killed the scalar. The deeper relational structure --- how a horse's effective
ability depends on the \emph{style composition} of \emph{this} field --- is the interacting theory
of Chapter~\ref{ch:theory}, not yet in production.
\end{remark}
